\section{Overview}
\label{sec:overview}

The Client Interface (CIF) accepts AXI read and write transactions, validates
them, divides each accepted transaction into row-bounded MC packet segments,
and returns AXI responses in order for each AXID.  Read and write traffic use
independent instances of the same transaction-tracking ROB architecture.

\subsection{Global assumptions}
\begin{itemize}[noitemsep]
  \item AXID is 6 bits (64 AXI IDs); each ROB instance has 64 per-AXID Slot
        FIFOs, each of depth four.
  \item One accepted AXI transaction produces 1--4 packet segments.  AXI
        beats are represented by a segment's \texttt{beat\_cnt}; a beat is
        not an MC packet.
  \item An AXI transaction must not cross 4~KB.  With runtime-configured
        \texttt{row\_size} of at least 1~KB, row crossings are at most three
        and \texttt{MAX\_SEGMENTS=4}.
  \item MC Core echoes the packet's \texttt{GLOBAL\_TAG} on completion.
        Completion-status encoding is \textbf{TBD with MC Core}.
\end{itemize}

\section{Parameters}
\begin{center}
\begin{tabular}{ll}
\toprule
\textbf{Parameter} & \textbf{Value}\\
\midrule
\texttt{AXID\_W} & 6\\
\texttt{ROB\_DEPTH} & 256 transactions per ROB instance\\
\texttt{ROB\_INDEX\_W} & 8\\
\texttt{PACKET\_NUMBER\_W} & 2\\
\texttt{GLOBAL\_TAG\_W} & 10: \texttt{\{ROB\_INDEX, PACKET\_NUMBER\}}\\
\texttt{MAX\_SEGMENTS} & 4\\
Per-AXID Slot FIFO & 64 FIFOs $\times$ depth 4, per ROB instance\\
Read Data Buffer watermark & 128 entries, owned independently by that buffer\\
\bottomrule
\end{tabular}
\end{center}

\section{Transaction identity and packet formation}
\label{sec:identity}

At AXI transaction acceptance, the applicable ROB Free List allocates
\texttt{ROB\_INDEX[7:0]}.  The index is pushed into the accepting AXID's Slot
FIFO only after the Burst Splitter has supplied
\texttt{splitter\_pkt\_count}.  The ROB initializes
\texttt{exp\_pkts=splitter\_pkt\_count} and clears its four-bit completion
bitmap before the entry becomes eligible for retirement.

At packet emission, the Burst Splitter assigns \texttt{PACKET\_NUMBER} from
its two-bit \texttt{seg\_idx} and Packet Formation attaches
\texttt{GLOBAL\_TAG=\{ROB\_INDEX,PACKET\_NUMBER\}}.  This is the complete
MC completion identity.

\section{Write path}
\label{sec:writepath}
\subsection{Request path}

The AW Metadata FIFO, W Data FIFO, Request Validator, Error Handler, Burst
Splitter, Packet Formation Unit, and Write Packet FIFO retain their existing
roles.  The accepting write ROB allocates the transaction index.  The splitter
generates one to four row-bounded segments and the Write Packet FIFO forwards
one packet per segment to MC Core.

\subsection{Write ROB and response}

The write ROB has its own 256-entry Free List, Metadata RAM, and 64 Slot
FIFOs.  Its Metadata RAM is directly indexed by \texttt{ROB\_INDEX} and
stores \texttt{exp\_pkts}, \texttt{completion\_bitmap[3:0]}, and aggregate
\texttt{resp\_status}.  MC completions set the bit selected by
\texttt{PACKET\_NUMBER} and update aggregate status.  The status mapping and
the worst-status precedence are \textbf{TBD}; no value is assumed here.

When the head transaction of an AXID Slot FIFO has all expected completion
bits set, retirement selects it in AXID order.  It holds BRESP valid until the
downstream AXI response path accepts it.  Only that acceptance pops the Slot
FIFO and returns \texttt{ROB\_INDEX} to the write Free List.  Write traffic
has no Read Data Buffer.

\section{Read path}
\label{sec:readpath}
\subsection{Request path}

The AR Metadata FIFO, Request Validator, Error Handler, Burst Splitter,
Packet Formation Unit, and Read Packet FIFO retain their existing roles.  The
read ROB allocates \texttt{ROB\_INDEX} at acceptance.  The splitter emits one
to four segments, each carrying a distinct \texttt{PACKET\_NUMBER}.

\subsection{Read ROB, Data Buffer, and serializer}

The independent read ROB has the same generic structures as the write ROB.
Its Metadata RAM additionally stores \texttt{dbuf\_ptr[0..3]}, one opaque
Read Data Buffer allocation handle per packet number.  At read packet issue,
the handle for that segment's allocation is recorded before the packet is
issued.  Completion data is routed through that handle according to the
MC/Data-Buffer interface contract, then
\texttt{completion\_bitmap[PACKET\_NUMBER]} is set.

Once a Slot FIFO head is complete, a Retiring Register holds its
\texttt{ROB\_INDEX} while the Read Serializer emits the associated buffer
allocations in packet-number order.  The MC/Data-Buffer interface contract
defines response-unit semantics, within-segment serialization, and allocation
release.  \texttt{RLAST} is asserted on the final AXI beat of the final
segment; only after that beat and every associated contract-defined release
may the Slot FIFO be popped and \texttt{ROB\_INDEX} returned to the read Free
List.  Read Data Buffer handshakes, latency, pointer representation, and MC
response-unit semantics remain \textbf{TBD}; its 128-entry watermark is its
own capacity control, not ROB occupancy control.

\section{Request and response packet formats}
\label{sec:pktformat}

This is the compact, standalone reference for the CIF-generated request
packet and the MC completion response, gathering fields already established
in \S\ref{sec:identity}, \S\ref{sec:splitter}, and \S\ref{sec:rob}.  The
request format (CIF\texorpdfstring{$\to$}{->}MC) and the response format
(MC\texorpdfstring{$\to$}{->}CIF) carry different degrees of freeze and must
not be conflated.

\subsection{CIF \texorpdfstring{$\rightarrow$}{->} MC Core request packet}

Packet Formation emits exactly one MC request packet per Segment Table entry
(\S\ref{sec:splitter}).  The packet is transferred to MC Core using the
\texttt{PKT\_VALID}/\texttt{PKT\_READY} handshake; that handshake is an
interface signal pair, not a field carried inside the packet.  Every request
packet, read or write, carries:

\begin{center}
\begin{tabular}{lll}
\toprule
\textbf{Field} & \textbf{Width} & \textbf{Meaning}\\
\midrule
\texttt{GLOBAL\_TAG} & 10 & \texttt{\{ROB\_INDEX[7:0],PACKET\_NUMBER[1:0]\}}; this segment's completion identity\\
\texttt{addr} & 32 & Segment start address\\
\texttt{beat\_cnt} & 8 & AXI beats represented by \emph{this segment}, not the transaction's packet count\\
packet type & inherited, unconfirmed width & Distinguishes Write/Read/Err\_write/Err\_read; set by the Request Validator/Error Handler and carried unchanged through the Burst Splitter to Packet Formation\\
\bottomrule
\end{tabular}
\end{center}

\textbf{Write request packets} additionally carry write data for the beats
represented by \texttt{beat\_cnt}.  The current documents establish only that
a W Data FIFO supplies this data on the write path (\S\ref{sec:writepath});
the per-segment write-data field width, byte-enable representation, and its
exact relationship to \texttt{beat\_cnt} are not fixed by any current CIF
document and remain \textbf{TBD}.

\textbf{Read request packets} carry no data field.  The interpretation of
\texttt{addr} and \texttt{beat\_cnt} at the MC request interface is governed
by the CIF--MC interface contract.

For WRAP bursts, the CIF-side WRAP Beat Address Table remains responsible for
generating the per-beat address sequence; the exact representation of that
sequence at the MC request interface is part of the CIF--MC interface
contract and is not fixed by this document.

Packet-type width/encoding is not restated by the post-refactor architecture
documents; the field's existence and the values it must distinguish (Write,
Read, Err\_write, Err\_read) are established, but its bit width should be
treated as inherited-but-unconfirmed, in the same manner as \texttt{status}
(see the Unresolved/TBD discussion referenced from \S\ref{sec:rob}).

\subsection{MC Core \texorpdfstring{$\rightarrow$}{->} CIF response interface}
\label{sec:respformat}

The current architecture requires only that a completion be attributable to
one segment and carry a status:

\begin{center}
\begin{tabular}{lll}
\toprule
\textbf{Field} & \textbf{Width} & \textbf{Meaning}\\
\midrule
\texttt{GLOBAL\_TAG} & 10 & Echoed back unmodified; decodes directly to \texttt{ROB\_INDEX} and \texttt{PACKET\_NUMBER}\\
\texttt{status} & TBD & Raw MC completion status; encoding \textbf{TBD}\\
\texttt{data} & TBD & Read response payload, read segments only\\
\bottomrule
\end{tabular}
\end{center}

The following remain \textbf{MC/Data-Buffer interface-contract dependent} and
are intentionally not fixed by this specification:
\begin{itemize}[noitemsep]
  \item response transport and handshake naming --- no \texttt{RESP\_VALID}/\texttt{RESP\_READY} signal pair is assumed;
  \item response payload width;
  \item response-unit granularity, i.e.\ whether one response unit corresponds to one AXI beat, one fixed-width data unit, or another unit;
  \item within-segment serialization semantics.
\end{itemize}

One segment's completion sets exactly one \texttt{completion\_bitmap} bit
regardless of that segment's \texttt{beat\_cnt} --- it is not one bit per AXI
beat.  A segment completing is therefore not the same event as \texttt{RLAST}
on the AXI R channel, which additionally requires that segment's beats to
have been serialized and accepted downstream.  See the read ROB's completion
and read-instance handling (\S\ref{sec:rob}) for how this response is
consumed once its transport is defined by that contract.

\section{Ordering, capacity, and reset}

Slot FIFO order is transaction acceptance order for each AXID and therefore
enforces AXI response ordering.  Metadata RAM capacity is transaction capacity
only.  A full Slot FIFO blocks allocation for that AXID; an empty Free List
blocks allocation globally.  The latter means all 256 transaction slots are
allocated, not merely that one Slot FIFO is full.  Read request issuance is
controlled separately by pooled Read Data Buffer availability.

Reset is active-low, asserted asynchronously, and deasserted synchronously;
\S\ref{sec:resetseq} gives the full initialization sequence.  It clears both
ROB Free Lists, Metadata RAM state, Slot FIFOs, and Retiring Registers.  The
Read Data Buffer is reset and drained according to its own specification.  A
soft reset is entered only after outstanding responses have been accepted and
both ROB instances have no allocated Slot FIFO entries.

\section{Error handling}

Invalid requests enter the normal packet path with the appropriate error
packet type.  They receive ordinary \texttt{GLOBAL\_TAG}s and complete through
the corresponding ROB.  Error status is preserved until ordered retirement;
the external mapping of raw MC status to AXI \texttt{BRESP/RRESP} remains
\textbf{TBD with MC Core}.
